% !TeX root = ../main.tex
\section{Характеристики шума на выходе ФСС: расчет номиналов элементов}

В таблице с вариантами лабораторной работы серым цветом выделен вариант №4. Для расчета будем использовать его исходные данные, а также общие параметры схемы:
\begin{itemize}
    \item Центральная частота (частота ПЧ): $f_{\text{п}} = 10 \text{ МГц} = 10^7 \text{ Гц}$
    \item Характеристическое сопротивление: $\rho = R_{\Gamma} = R_H = 1.0 \text{ кОм} = 10^3 \text{ Ом}$
    \item Полоса пропускания ФСС: $\Pi_{\text{ФСС}} = \Pi_{\text{УПЧ}} = 1.0 \text{ МГц} = 10^6 \text{ Гц}$
\end{itemize}

Согласно заданию, рассчитанные номиналы должны иметь точность не менее 3-х знаков после запятой для предотвращения искажения АЧХ.

\vspace{1em}
\textbf{1. Расчет емкости связи ($C_{\text{св}}$)}

Используем формулу $C_{\text{св}} = \frac{1}{2\pi f_{\text{п}} \rho}$:
\begin{align*}
    C_{\text{св}} &= \frac{1}{2 \cdot \pi \cdot 10^7 \cdot 10^3} = \frac{1}{2\pi \cdot 10^{10}} \approx 15.915 \cdot 10^{-12} \text{ Ф} = 15.915 \text{ пФ}
\end{align*}

\vspace{1em}
\textbf{2. Расчет емкости внутреннего контура ($C_2$)}

Используем формулу $C_2 = \frac{1}{\pi \Pi_{\text{ФСС}} \rho} - 2C_{\text{св}}$:
\begin{align*}
    C_2 &= \frac{1}{\pi \cdot 10^6 \cdot 10^3} - 2 \cdot 15.915 \cdot 10^{-12} \\
    C_2 &= \frac{1}{\pi \cdot 10^9} - 31.830 \cdot 10^{-12} \\
    C_2 &\approx 318.310 \cdot 10^{-12} - 31.830 \cdot 10^{-12} = 286.480 \cdot 10^{-12} \text{ Ф} = 286.480 \text{ пФ}
\end{align*}

\vspace{1em}
\textbf{3. Расчет индуктивности внутреннего контура ($L_2$)}

Используем формулу $L_2 = \frac{\Pi_{\text{ФСС}} \rho}{4\pi f_{\text{п}}^2}$:
\begin{align*}
    L_2 &= \frac{10^6 \cdot 10^3}{4\pi \cdot (10^7)^2} = \frac{10^9}{4\pi \cdot 10^{14}} = \frac{10^{-5}}{4\pi} = \frac{10}{4\pi} \cdot 10^{-6} \approx 0.796 \cdot 10^{-6} \text{ Гн} = 0.796 \text{ мкГн}
\end{align*}

\vspace{1em}
\textbf{4. Расчет емкости крайних контуров ($C_1$)}

Используем формулу $C_1 = C_2 / 2$:
\begin{align*}
    C_1 &= \frac{286.480 \text{ пФ}}{2} = 143.240 \text{ пФ}
\end{align*}

\vspace{1em}
\textbf{5. Расчет индуктивности крайних контуров ($L_1$)}

Используем формулу $L_1 = 2L_2$:
\begin{align*}
    L_1 &= 2 \cdot 0.796 \text{ мкГн} = 1.592 \text{ мкГн}
\end{align*}

\vspace{1em}
\textbf{Итоговые номиналы для моделирования (вариант 4):}
\begin{itemize}
    \item $C_{\text{св}} = 15.915 \text{ пФ}$
    \item $C_1 = 143.240 \text{ пФ}$
    \item $L_1 = 1.592 \text{ мкГн}$
    \item $C_2 = 286.480 \text{ пФ}$
    \item $L_2 = 0.796 \text{ мкГн}$
\end{itemize}